% !TEX root = ../../main.tex

\section{系统技术路线概述}

HearBridge 系统面向听障人士交互训练与辅助交流场景，整体采用多端协同、前后端分离和识别服务独立部署的技术路线。
系统由 HarmonyOS 移动端、Spring Boot 服务端、Vue 后台管理端和 Python 手势识别服务组成。
移动端负责用户侧的手语资源浏览、动作示范、训练交互、实时识别、视频选择和识别结果展示；
服务端负责用户认证、业务数据管理、文件资源管理、识别服务调用和语义修正结果组织；
后台管理端负责手势分类、手势资源、样本数据、训练任务和模型版本管理；
Python 手势识别服务负责关键点提取、特征转换、模型训练、实时识别和句子视频识别。

从系统实现角度看，该技术路线能够将不同类型的任务分配给更适合的技术环境。
HarmonyOS 端适合承载移动端交互、多媒体选择和训练结果展示；
Spring Boot 适合完成业务接口、权限校验和数据管理；
Vue 后台管理端适合构建表格、表单、任务触发和资源维护等管理页面；
Python 生态则更适合集成 MediaPipe、TensorFlow/Keras、OpenCV 等视觉识别和模型训练工具。
通过这种分层协同方式，系统能够在保证移动端交互体验的同时，使识别模型训练、资源管理和模型版本发布具备较好的维护能力。

本章主要介绍系统实现过程中涉及的关键技术，包括 HarmonyOS 客户端技术、服务端技术、后台管理端技术、手势识别服务技术、句子视频识别技术以及接口调试与联调支撑技术。
具体的系统需求、结构设计、核心实现和实验结果将在后续章节中进一步展开。
